\documentclass[11pt]{article}
\usepackage[margin=1in]{geometry}
\usepackage[T1]{fontenc}
\usepackage[utf8]{inputenc}
\usepackage{array}
\usepackage{longtable}
\usepackage{booktabs}
\usepackage{hyperref}

\hypersetup{colorlinks=true, linkcolor=blue, urlcolor=blue}

\newcommand{\src}[1]{\texttt{\detokenize{#1}}}

\title{Sections 3--4 Formalization Overview}
\author{GeoTop formalization notes}
\date{2026-06-17}

\begin{document}
\maketitle

\section{Current Status}

The Moise Sections 3--4 theorem chain is closed in the current \src{dev34}
split.  Section 3 is represented by \src{Theorem_GT_3_1} through
\src{Theorem_GT_3_7}; Section 4 is represented by \src{Theorem_GT_4_1}
through \src{Theorem_GT_4_10}.  The main definitions are
\src{geotop_free_2_simplex}, \src{geotop_boundary_free_2_simplex}, and
\src{geotop_brick_decomposition}.  The supporting invariance-of-domain theorem
is \src{Theorem_GT_4_invariance_of_domain} in \src{gb/GeoTopBase.thy:6722}.

A direct scan of the current Section 3--4 theory files found no active
\src{sorry} or \src{oops} markers.  The only hit for \src{admit} was prose in a
comment using the English word ``admits''.

\section{Size and Location}

\begin{longtable}{p{0.18\textwidth}p{0.35\textwidth}p{0.16\textwidth}p{0.23\textwidth}}
\toprule
\textbf{Area} & \textbf{File or directory} & \textbf{Size} & \textbf{Role}\\
\midrule
\endhead
Base prefix & \src{dev34_prefix_base/GeoTop_3_4_Prefix_Base.thy} & 5,865 lines & Theorems 3.1--3.2 and simplex/homeomorphism setup.\\
Mid prefix & \src{dev34_prefix_mid/GeoTop_3_4_Prefix_Mid.thy} & 67,258 lines & Free-simplex machinery, Theorems 3.3--3.7, and Theorems 4.1--4.2.\\
Prefix & \src{dev34_prefix/GeoTop_3_4_Prefix.thy} & 32,860 lines & Theorem 4.3, brick-decomposition, and Theorem 4.4.\\
Facts layer & \src{dev34_facts/GeoTop_3_4_Facts.thy} & 7,516 lines & Theorems 4.5--4.7 and reusable Jordan/separation facts.\\
Final layer & \src{dev34/GeoTop_3_4.thy} & 32,239 lines & Theorems 4.8--4.10 and manifold/link/star work.\\
Main proof files & above five files & 145,738 lines & Working Section 3--4 formalization split.\\
Fast cache & \src{.dev34_fast_cache/} & about 90 MB & Generated prefix/slice cache for focused checks.\\
Report snippets & \src{sections_3_4_snippets/} & about 160 KB & Cached exact source excerpts for the detailed report.\\
Rendered snippets & \src{sections_3_4_rendered/} & about 80 KB & Rendered book excerpts for the detailed report.\\
\bottomrule
\end{longtable}

\section{How It Was Done}

The development was split into layers so finished infrastructure could be
reused without replaying the whole theory stack on every edit.  The main layers
are \src{dev34_prefix_base}, \src{dev34_prefix_mid}, \src{dev34_prefix},
\src{dev34_facts}, and \src{dev34}.

The workflow relied heavily on searchable indexes: \src{THEOREMS_AND_DEFS.txt}
for declaration names and line numbers, \src{STMT_INDEX.txt} for normalized
statements, and full-tree \src{rg} searches for local proof facts and comments.
The fast iteration path used \src{.dev34_fast_cache} and
\src{check_dev34_fast.sh}.  The cache stores generated prefixes and focused
slices so that long theorem packages can be checked without rebuilding all
earlier finished text.  It is a development accelerator, not a substitute for
final Isabelle verification.

\section{Proof Strategy}

Section 3 was closed by turning Moise's geometric induction into explicit
reusable packages.  Theorem 3.3 uses the stronger ``at least two free
2-simplexes'' induction.  Theorem 3.4 uses the Figure 3.3 fold and inverse-fold
packages.  Theorem 3.7 reuses the same fold normalization with support control
outside a chosen open set.

Section 4 combines faithful formalizations of Moise's arguments with
library-backed topological bridges.  Theorem 4.2 packages the opposite-boundary
arc separation argument.  Theorem 4.4 packages the brick/regular-neighborhood
component-transfer argument.  Theorems 4.3, 4.5, and 4.6 use
HOL-Analysis/Jordan-Brouwer bridges where that is cleaner than replaying every
planar picture argument.  Theorem 4.8 follows Moise's five-link-lemma plan
inside the theorem body.  Theorem 4.9 identifies the topological boundary with
one-incident edges, and Theorem 4.10 uses invariance of domain for the
frontier-to-boundary inclusion.

\section{Problems and Solutions}

\begin{longtable}{p{0.43\textwidth}p{0.50\textwidth}}
\toprule
\textbf{Problem} & \textbf{Solution}\\
\midrule
\endhead
Long theory files made ordinary verification slow. & Split the development into layers and used \src{.dev34_fast_cache} focused prefix/slice checks.\\
Search by memory was unreliable in a 100k+ line development. & Regenerated and used \src{THEOREMS_AND_DEFS.txt}, \src{STMT_INDEX.txt}, and full-index \src{rg} searches frequently.\\
Moise's proofs rely on figures and phrases such as ``evidently''. & Replaced picture arguments by named local lemmas and reusable geometric packages.\\
Theorem 3.3 needed a stronger induction than the final theorem statement. & Proved and used the two-free-simplex count theorem, then extracted the weaker witness theorem.\\
Theorems 3.4 and 3.7 required controlled plane homeomorphisms. & Built supported Figure 3.3 fold/inverse-fold packages and a supported normalization theorem.\\
Theorem 4.4's brick proof was too broad if formalized as a general brick theory. & Focused on the exact component-transfer conclusion needed by the theorem.\\
Theorem 4.8 needed many local incidence and link facts. & Broke the proof into L1--L5 local steps plus named semicircle, three-incident-simplex, and Figure 4.10 cone/link packages.\\
The detailed LaTeX report became slow to compile when it read high line ranges from large \src{.thy} files. & Cached source excerpts in \src{sections_3_4_snippets/} and rendered book excerpts in \src{sections_3_4_rendered/}.\\
\bottomrule
\end{longtable}

\section{Outcome and Caveats}

The split architecture made it possible to close very large theorem packages
incrementally.  The indexes made theorem lookup and line-number reporting
practical, and the proof packages now document the correspondence between
Moise's informal geometric steps and formal Isabelle objects.

The main caveats are structural.  The \src{dev34} split is a working
formalization layer; integration into the final canonical \src{GeoTop.thy}
layout may still require cleanup.  Some book sublemmas are absorbed into
stronger library-backed theorems rather than reproduced line-for-line.  The
fast cache is generated local state and should be treated as an iteration aid.
The local TeX installation in this environment cannot build PDFs because
required TeX format files are missing, but the LaTeX source is intended to be
portable.

\end{document}

